\section{Подсистема телеметрии и интеграция с внешними клиентами}

Сбор и трансляция метрик загруженности дорог в реальном времени (Heatmap) является ресурсоемкой задачей, способной существенно замедлить работу основного вычислительного конвейера. Для решения этой проблемы в \texttt{traffic\_core} реализован асинхронный модуль телеметрии, интегрированный с шиной сообщений ZeroMQ.

\subsection{Асинхронная архитектура трансляции состояния (ZeroMQ)}

Вместо прямых TCP-соединений или HTTP-опросов со стороны внешних систем, ядро симуляции выступает в роли однонаправленного источника данных (Publisher). Для организации потоковой передачи применена легковесная библиотека ZeroMQ (паттерн Pub-Sub). 

На этапе «Шаг телеметрии» в главном цикле (Main Loop) симулятор (\texttt{data\_provider}) формирует массив только измененных состояний сегментов (дельт). Если плотность потока на ребре изменилась, формируется компактная структура данных, содержащая \texttt{edge\_id} и текущий расчетный штраф BPR (или цветовую кодировку загруженности). 
Сформированные массивы бинарных данных сериализуются (с применением структур, аналогичных Protobuf/FlatBuffers для обеспечения компактности) и атомарно передаются в фоновый поток Publisher-а ZeroMQ через неблокирующую очередь \texttt{ConcurrentQueue}. 

Такая архитектура гарантирует, что если внешний клиент (например, API Gateway или Web-браузер) отключится, работает медленно или не успевает читать сокет, внутренний цикл సిмулятора C++ не заблокируется (Non-blocking I/O). Сообщения накапливаются во внутренних буферах ZeroMQ или сбрасываются при переполнении, сохраняя строгий бюджет времени на симуляцию.

\subsection{Формирование Heatmap и агрегация метрик}

Для визуализации дорожной обстановки (тепловой карты заторов — Heatmap) на клиентской стороне, \texttt{traffic\_core} не занимается рендерингом тайлов. Вычислительное ядро транслирует только сырые факты: идентификатор направленного ребра и уровень его перенасыщения (отношение $V(t) / C_{max}$). 

МПР (Модуль принятия решений) также использует этот телеметрический контур для сбора глобального показателя эффективности системы — Travel Time Index (TTI). Поскольку ядро (\texttt{router}) фиксирует плановое время прибытия (ETA), а \texttt{data\_provider} отслеживает фактическое завершение поездки агента, система на лету агрегирует показатель $TTI = W_{real} / T_{free}$ для каждого успешно доехавшего агента. 

Асинхронный экспорт этих метрик через ZeroMQ позволяет внешнему API Gateway (написанному на Python или Go) собирать статистику, сохранять исторические срезы в базу данных (ClickHouse/PostgreSQL) и раздавать агрегированные тепловые карты десяткам веб-клиентов (на базе MapLibre GL JS) по WebSocket без какого-либо дополнительного влияния на ядро маршрутизатора.
